\section{Decomposition at the Municipality Level}
\label{sec:decomposition_mun}

The decomposition in Section~\ref{sec:decomposition} produces national-level bounds on exit and re-labeling using event-time coefficients estimated across treated municipalities. Those coefficients are averages: they tell us how the typical treated municipality's registry composition shifts at adoption, but not how exit and re-labeling distribute across municipalities or across age groups within them. For the welfare calibration in Section~\ref{sec:calibration} and for the spatial visualization in Section~\ref{ssec:decomposition_spatial}, a municipality-level decomposition is necessary. This section develops one.

The municipality-level framework cannot rely on the never-treated comparison group that anchors the cross-municipality identification. Instead, it uses a different identifying restriction: voters age mechanically. A voter who is 30 at one election is 32 at the next, regardless of whether their municipality has adopted BVR. The mechanical aging structure, combined with municipality-specific pre-treatment survival rates, generates a counterfactual prediction of what each age cohort's count should look like at the post-adoption election absent BVR. The gap between predicted and actual cohort counts identifies BVR-attributable losses. Because the framework requires age-disaggregated counts, it relies on TSE's age-by-education panel, which is available from 2008 onward. The 2008 baseline replaces the 2006 baseline used in the cross-section decomposition.

Two design choices distinguish the municipality-level framework from its national counterpart. First, identification comes from within-municipality variation across time rather than from the contrast between treated and never-treated municipalities. Second, the framework explicitly disaggregates the registry into joint age-by-education cells, which allows the cleanest re-labeling identification in age cohorts past the typical schooling-completion window. The cost is that the framework conflates BVR-induced exit with voluntary out-migration, which the cross-section design absorbs through the never-treated comparison. Section~\ref{ssec:decomposition_mun_math} develops the framework. Section~\ref{ssec:decomposition_mun_survival} discusses estimation of the natural survival rates that anchor the counterfactual. Section~\ref{ssec:decomposition_mun_relabel} describes the re-labeling identification within age cohorts. Section~\ref{ssec:decomposition_mun_interp} discusses interpretation and limitations, with particular attention to out-migration.

\subsection{The Cohort-Aging Framework}
\label{ssec:decomposition_mun_math}

For municipality $m$ at election year $t$, let $N_{m,t,c,e}$ denote the count of registered voters in age cohort $c$ and education category $e$. The cohort index $c$ runs over the age bands reported by TSE, which combine single-year bins for the youngest registration ages with five-year bins for older voters and an absorbing top bin.\footnote{The TSE reports 22 age bins in total: single-year bins for ages 16, 17, 18, 19, and 20; a four-year bin for ages 21-24; fifteen five-year bins from 25-29 through 95-99; and the absorbing 100+ category. Voters with invalid age records (TSE label \emph{Inv\'alido}) are excluded from the framework and reported separately as a small residual. The 21-24 bin spans four years of age while subsequent five-year bins span five; the mechanical-aging operator described below accounts for this asymmetry. The unknown-education category is small (about 0.08\% of the electorate) and is treated as in Section~\ref{sec:decomposition}: excluded from the framework. The age-by-education panel jointly observed at the municipality level starts in 2008; earlier TSE files report only the marginal education distribution.} The education index $e \in \{L, H\}$ takes the low- and high-education values from Section~\ref{sec:decomposition}. Brazilian elections occur on a two-year cycle, so the natural unit of analysis is the increment $t \to t+2$, corresponding to two years of voter aging within each age cohort.

The mechanical aging restriction says that, absent any net migration or BVR-induced cancellation, the count of voters in cohort $c$ at election $t+2$ equals the surviving share of voters who were in cohort $c$ or its predecessor at election $t$, plus the inflow of newly-eligible voters into the youngest cohort. Formally, define the mechanical-aging operator $\mathcal{T}$ acting on a vector of cohort counts. The operator distributes voters from cohort $c$ at $t$ into cohorts at $t+2$ according to where their ages fall after two years.

For the single-year bins, the entire cohort moves forward exactly two cells over a two-year cycle.\footnote{Voters aged 16 at $t$ are aged 18 at $t+2$, voters aged 17 at $t$ are aged 19 at $t+2$, voters aged 18 at $t$ are aged 20 at $t+2$, and voters aged 19 or 20 at $t$ enter the next bin (21-24) at $t+2$.} For the four-year bin, $\tfrac{2}{4}$ of voters cross into the next bin while $\tfrac{2}{4}$ remain. For each five-year bin, $\tfrac{2}{5}$ of voters move forward and $\tfrac{3}{5}$ remain. The 100+ bin is absorbing: all voters who reach it stay there. Under a uniform within-cohort age distribution, these fractions describe the proportion of each cohort's age range that crosses the next boundary over the two-year cycle.

Let $\sigma_{m,c}$ denote the natural survival rate of voters in cohort $c$ in municipality $m$ between two consecutive elections. The rate $\sigma_{m,c}$ captures all baseline processes that reduce the count: mortality and voluntary deactivation. It is bounded between zero and one and is typically close to one for younger cohorts and lower for older cohorts due to mortality. Section~\ref{ssec:decomposition_mun_survival} discusses estimation of $\sigma_{m,c}$ from pre-treatment data.

Let $I_{m,t+2}$ denote the inflow of newly-eligible voters into the age-16 bin at $t+2$. Voters become eligible to vote at 16 and required to vote at 18, so the age-16 bin at $t+2$ consists of newly-registered 16-year-olds; the age-17 bin at $t+2$ consists of those who were age 16 voters at $t$ who survived the cycle plus newly-registered 17-year-olds. Inflows continue at each younger age until the 18-year-old bin, where compulsory registration kicks in and inflows become structurally larger. The framework absorbs these inflows in the counterfactual prediction for the youngest cohorts using municipality-specific pre-treatment inflow rates.

The counterfactual prediction of cohort $c$'s count at $t+2$ absent BVR is
\begin{equation}
\widehat{N}_{m,t+2,c}^{\text{no-BVR}} = \sigma_{m,c'} \cdot \left[\mathcal{T}(N_{m,t,\cdot})\right]_c + I_{m,t+2,c},
\label{eq:counterfactual}
\end{equation}
where $\left[\mathcal{T}(N_{m,t,\cdot})\right]_c$ denotes the entry of the aged-forward vector corresponding to cohort $c$, $c'$ is the source cohort that contributed to $c$, and $I_{m,t+2,c}$ is the inflow of newly-registered voters into cohort $c$ at $t+2$. Inflows are nonzero only for the youngest age bins (16, 17, 18, 19); for cohorts 20 and older, $I_{m,t+2,c} = 0$ since voters cannot enter these bins except by aging up from younger bins. For interior cohorts, equation~\eqref{eq:counterfactual} weights the staying mass and the inflow from the immediately younger cohort by the appropriate cohort-specific survival rate.\footnote{The framework treats survival rates as cohort-specific within municipality but does not explicitly model how survival differs between cells with different education levels within the same age cohort. This is a simplifying assumption; in practice, mortality differentials by education within age cohorts may be modest at most ages and become noticeable only at advanced ages.}

The actual count $N_{m,t+2,c}$ differs from the no-BVR prediction $\widehat{N}_{m,t+2,c}^{\text{no-BVR}}$ when BVR-induced exit removes voters from cohort $c$. Define BVR-induced exit by cohort as
\begin{equation}
E_{m,t+2,c} = \max\left\{0, \widehat{N}_{m,t+2,c}^{\text{no-BVR}} - N_{m,t+2,c}\right\}.
\label{eq:exit_cohort}
\end{equation}
The non-negativity floor reflects the substantive constraint that BVR cannot create voters; if observed exceeds predicted, the gap is attributed to estimation noise rather than negative exit. The total municipality-level exit at the post-adoption election is the sum across cohorts:
\begin{equation}
E_{m,t+2} = \sum_c E_{m,t+2,c}.
\label{eq:exit_total}
\end{equation}

\subsection{Estimating the Natural Survival Rates}
\label{ssec:decomposition_mun_survival}

The cohort-specific survival rate $\sigma_{m,c}$ is the parameter that anchors the counterfactual prediction. Because it is not directly observable, it must be estimated. The framework offers three approaches, listed in increasing order of municipality specificity.

\textbf{National survival rates from never-treated municipalities.} The simplest approach uses never-treated municipalities to estimate population-weighted average survival rates by cohort. For each pair of consecutive election years where neither year reflects BVR exposure, compute the realized cohort-to-cohort transition implied by the data and back out the implied survival rate after netting out the mechanical aging operator. This produces a single $\sigma_c$ per cohort, applied uniformly across municipalities.

\textbf{State-level survival rates.} A middle-ground approach estimates $\sigma_{c,s}$ separately for each state $s$, using never-treated municipalities or pre-treatment cycles within the state. This allows for state-specific baseline mortality and migration patterns while pooling enough observations to keep the estimates stable.

\textbf{Municipality-specific pre-treatment survival rates.} The most municipality-specific approach uses the within-municipality transitions at $t \to t+2$ for years $t$ that pre-date the municipality's BVR adoption. For municipalities first treated in 2014, the cycles 2008-2010 and 2010-2012 provide two pre-treatment transitions from which to estimate $\hat\sigma_{m,c}$. For 2016 and 2018 cohorts, more pre-treatment cycles are available. This approach captures any unobserved municipality-specific factors that affect the survival rate, at the cost of noisier estimates in small municipalities.

The municipality-specific approach is the most defensible in principle because it absorbs unobserved heterogeneity in mortality, migration patterns, and registration practices. The trade-off is statistical precision: with 22 cohorts per municipality-year, the per-cohort population is much smaller than the cross-section decomposition's two education categories, and in small municipalities a single cohort cell may contain only a few dozen voters. The realized two-cycle survival rate has substantial sampling noise in such cells, and the resulting exit estimate inherits that noise. The framework as implemented uses the municipality-specific estimate when the cohort population exceeds a stability threshold (typically 100 voters for the youngest single-year bins and 200 voters for the five-year bins) and falls back to the state-level estimate otherwise. As a further robustness check, the framework can aggregate the 22 fine cohorts into coarser bands (for instance, 16-19, 20-24, 25-34, 35-44, 45-54, 55-64, 65+) before computing survival rates; the coarser aggregation produces less precise heterogeneity but more stable per-municipality estimates.

The pre-treatment estimation step is precisely where the parallel-trends idea from the cross-section decomposition appears in the municipality-level framework. In the cross-section, parallel trends justify the never-treated comparison; here, the assumption is that municipality $m$'s pre-treatment cohort transitions are informative about its post-treatment counterfactual transitions. This is a within-municipality version of the same identifying logic. It is testable indirectly by examining the stability of pre-treatment transitions across multiple cycles in the same municipality. Wide swings in pre-treatment transitions suggest the municipality has unstable baseline dynamics and the framework's prediction will be noisy.

\subsection{Re-labeling Identification within Age Cohorts}
\label{ssec:decomposition_mun_relabel}

The cohort-aging framework identifies total exit by age cohort, but it does not directly identify re-labeling, since re-labeling shifts voters from low-education to high-education within the same age cohort without changing cohort membership. To recover re-labeling, the framework looks at the education composition within each age cohort at $t+2$ relative to its predecessor at $t$.

For each cohort $c$ at time $t+2$, define the predicted education composition under the no-BVR counterfactual:
\begin{equation}
\widehat{N}_{m,t+2,c,e}^{\text{no-BVR}} = \widehat{N}_{m,t+2,c}^{\text{no-BVR}} \cdot \pi_{m,t,c,e} + \delta_{c,e} \cdot \widehat{N}_{m,t+2,c}^{\text{no-BVR}},
\label{eq:counterfactual_edu}
\end{equation}
where $\pi_{m,t,c,e}$ is municipality $m$'s share of cohort $c$ in education category $e$ at the prior election, and $\delta_{c,e}$ is a cohort-specific natural education progression rate. The first term carries forward the prior education composition; the second term adjusts for the share of cohort-$c$ voters who naturally complete additional schooling between elections.

For older cohorts, $\delta_{c,e}$ is approximately zero. Voters past 35 rarely complete formal schooling at meaningful rates over a two-year window, and the rate falls further for cohorts 45-49 and older. For these cohorts, equation~\eqref{eq:counterfactual_edu} simplifies to a pure carry-forward: the predicted education composition equals the prior education composition. Any observed shift from low-education to high-education within these cohorts at $t+2$ is then attributable to either re-labeling or to differential exit by education within the cohort. With the cohort-level total exit pinned down by equation~\eqref{eq:exit_cohort}, allocating it across education categories under a no-high-education-exit assumption (analogous to Scenario B in Section~\ref{sec:decomposition}) recovers a point estimate for re-labeling within the cohort:
\begin{equation}
R_{m,t+2,c} = \widehat{N}_{m,t+2,c,L}^{\text{no-BVR}} - N_{m,t+2,c,L} - E_{m,t+2,c},
\label{eq:relabel_cohort}
\end{equation}
where $E_{m,t+2,c}$ is allocated entirely to the low-education category under Scenario B. The total municipality-level re-labeling at the post-adoption election is the sum across cohorts:
\begin{equation}
R_{m,t+2} = \sum_c R_{m,t+2,c}.
\label{eq:relabel_total}
\end{equation}

For younger cohorts, $\delta_{c,e}$ must be estimated rather than assumed to be zero. Voters in the 16-24 range, in particular, are likely to complete additional schooling between elections: those completing secondary school typically do so between 17 and 19, and those completing higher education typically do so between 22 and 24. The natural education progression rate for these cohorts is identified from pre-treatment within-municipality transitions, in the same way as the survival rate. The cleaner re-labeling identification therefore concentrates on the cohorts aged 35 and older,\footnote{Specifically the TSE bins 35-39, 40-44, 45-49, 50-54, 55-59, 60-64, 65-69, 70-74, 75-79, 80-84, 85-89, 90-94, 95-99, and the absorbing 100+ category.} where natural progression is small and the identifying assumption is most credible. The 25-29 and 30-34 cohorts occupy an intermediate position where some completion of higher education or graduate study still occurs, but at lower rates than the youngest bins.

A further refinement uses the cell-level biometric uptake to bound re-labeling from above. Re-labeling requires recadastramento: a voter's education record is updated only when the voter physically attends the cartório eleitoral and completes the biometric registration process. Voters who have not undergone recadastramento retain their original education record by construction. Let $b_{m,t,c}$ denote the share of cohort $c$ in municipality $m$ at year $t$ with completed biometric registration, observed directly in the data. Then the re-labeling share within cohort $c$ at $t+2$ is bounded above by
\begin{equation}
R_{m,t+2,c} \leq b_{m,t+2,c} \cdot N_{m,t+2,c}.
\label{eq:relabel_pctbvr_bound}
\end{equation}
This bound is mechanical: a voter cannot be re-labeled without having been through recadastramento. In strict municipalities post-treatment, $b_{m,t+2,c}$ approaches one for all cohorts, so the bound is non-binding and the re-labeling estimate from equation~\eqref{eq:relabel_cohort} stands. In hybrid municipalities, $b_{m,t+2,c}$ is meaningfully below one (the descriptive evidence in the rho estimation analysis showed average voluntary uptake of 19.5\% at adoption), so the bound tightens the re-labeling estimate substantially. The implementation reports both the framework's main estimate and the $b$-bounded estimate; cells where the bound binds are flagged separately, since they indicate either small-sample noise in the cohort-aging counterfactual or model misspecification.

\subsection{Interpretation and Limitations}
\label{ssec:decomposition_mun_interp}

The municipality-level decomposition produces, for each treated municipality at its post-adoption election, an estimated number of BVR-induced exits and re-labelings, decomposed by age cohort. These estimates feed the spatial visualization in Section~\ref{ssec:decomposition_spatial} and the welfare calibration in Section~\ref{sec:calibration}. They also support a heterogeneity narrative that the cross-section decomposition cannot: where exit and re-labeling are concentrated within the population, by age and by joint age-education category.

Two interpretive points warrant emphasis before discussing limitations. First, the framework attributes the gap between predicted and actual cohort counts to BVR-induced exit, but the gap could in principle reflect other shocks specific to the municipality between $t$ and $t+2$. The framework is therefore best interpreted as identifying the BVR-induced loss \emph{conditional on the absence of other large idiosyncratic shocks}; absent such shocks, the framework gives a clean within-municipality identification. Second, the no-high-education-exit assumption used to allocate cohort exit across education categories within Scenario B is the same as in the cross-section decomposition. The Scenario C alternative (proportional exit) can be applied at the cohort level by replacing the Scenario B allocation with a proportional one, producing a corresponding range of re-labeling estimates per municipality.

The most important limitation is the conflation of BVR-induced exit with voluntary out-migration. Brazilian voters are required to keep their registration current to their domicile, which means voters who relocate from one municipality to another transfer their registration. From the perspective of the source municipality, an out-migrant looks identical to a BVR cancellation: the voter disappears from the registry. The framework cannot separate these two channels.

This is a substantive concern. Brazil has high inter-municipal mobility, particularly in metropolitan regions and in municipalities affected by economic shocks. Over a two-year window, several percent of registered voters in any given municipality may have transferred their registration for non-BVR reasons. The cross-section decomposition handles this through the never-treated comparison: voluntary migration affects all municipalities, so the differential between treated and never-treated isolates BVR. The municipality-level framework loses this protection. Its exit estimate therefore captures BVR-induced cancellation \emph{plus} any deviation in out-migration during the post-adoption period from the municipality's pre-treatment migration baseline. If the pre-treatment survival rate $\hat\sigma_{m,c}$ already reflects the municipality's typical out-migration pattern, the framework subtracts that baseline and what remains is plausibly BVR-attributable. If migration patterns shift between the pre-treatment and post-adoption periods for reasons unrelated to BVR, the framework misattributes some of the migration shift to BVR.

Two implications follow. First, the municipality-level exit estimate should be interpreted as an upper bound on BVR-induced exit, not as a point estimate. The estimate captures BVR cancellation plus any post-adoption out-migration shock unrelated to BVR. In municipalities with stable migration patterns, the upper bound is close to the true BVR effect; in municipalities with substantial migration shifts, the gap is larger. Second, the cross-section decomposition in Section~\ref{sec:decomposition} remains the cleaner identification of BVR's causal effect on the registry. The municipality-level framework supplements but does not supersede the cross-section results. The two frameworks serve different purposes: the cross-section identifies the average causal effect of BVR with the strongest available identification; the municipality-level decomposition produces spatial heterogeneity in the exit and re-labeling estimates under a within-municipality identifying restriction that is less defensible but operationally necessary for the spatial application.

A second limitation concerns survival rate estimation in small municipalities. Brazilian municipalities range from a few hundred voters to several million; survival rate estimates are noisy in the smallest municipalities. The framework's fallback to state-level survival rates for small municipalities mitigates this but does not eliminate it. Municipality-level results in the smallest decile of the size distribution should be treated with caution.

A third limitation concerns the natural education progression rate for the youngest cohorts. Voters in the 16-24 age range do complete additional schooling between elections, and the empirical pre-treatment estimate of $\delta_{c,e}$ for these cohorts may itself be contaminated by recadastramento occurring outside formal BVR (during routine registration updates that occur for reasons such as address changes or the obligatory transition from the optional 16-17 voting status to the compulsory 18+ status). The framework's re-labeling identification is therefore cleanest in the cohorts aged 35 and older, where natural progression is small. Younger-cohort estimates should be reported but interpreted with this caveat.

A fourth limitation concerns the uniform within-cohort age distribution implicit in the mechanical-aging operator. The framework treats all voters within a five-year cohort as equally likely to cross into the next cohort over a two-year cycle, which is exactly correct only if the within-cohort age distribution is uniform across single-year ages. In practice, the within-cohort distribution is approximately uniform for the working-age cohorts where mortality is low, but increasingly tilted toward the younger end of the cohort for older bins where differential mortality thins the older sub-ages over time. The implied error in the aging operator is small for the cohorts that matter most for re-labeling identification (those between approximately 35 and 65, where mortality is low) and grows in the older cohorts (70 and above) where it interacts with already-noisy survival rate estimates. Census data on single-year-of-age population shares could in principle replace the uniform assumption with empirically observed within-cohort distributions, providing more precise mechanical-aging weights at the cost of additional implementation complexity. We do not pursue this refinement because the marginal precision gain is modest in the cohort ranges most relevant for the framework's main identification, and because the older cohorts where the assumption is most problematic are also those where survival rates are noisiest for reasons unrelated to the within-bin distribution. As a robustness check, the framework can be re-implemented under an alternative assumption (for instance, an age-tilted within-cohort distribution that approximates differential mortality) to verify that results are not driven by the uniformity assumption.

The core conclusion of the municipality-level decomposition is that, despite the limitations, the framework produces internally consistent estimates of exit and re-labeling at the municipality level. Comparison against the cross-section national totals from Section~\ref{sec:decomposition} requires care because the two frameworks use different baselines: the cross-section uses the 2006 registered electorate as its denominator, while the municipality-level framework uses 2008 because the age-by-education panel is unavailable for earlier years. A direct comparison of aggregated municipality-level estimates against the cross-section coefficients $\gamma_j^k$ is therefore not exact; it requires either rescaling both frameworks to a common baseline or interpreting the comparison qualitatively. We report the comparison qualitatively: aggregated municipality-level estimates and cross-section totals should agree in direction and rough magnitude, with deviations attributable to the baseline difference, to out-migration noise, and to small-municipality survival rate noise. Substantial qualitative divergence between the two approaches---for example, the municipality-level framework producing exit estimates several times larger than the cross-section, or estimates with systematically different signs---would indicate that the within-municipality identifying assumptions are violated, and the framework should be re-examined accordingly. The implementation reports the consistency check explicitly and flags any cohorts or sub-samples where the divergence is large.
